% !TeX program = xelatex
% ============================================================
% 系统开发工具基础 · 实验报告（第三周）
% 学号 25090032014  甘如嫣
%
% 注意：老师要求「每个人设计自己的模板」。
%       下面这套排版（配色 / 页眉页脚 / 代码样式）可自行调整，
%       换成你自己喜欢的配色和字体，它才真正算"你的模板"。
% ============================================================
\documentclass[12pt,a4paper]{ctexart}

% ---------- 页面与包 ----------
\usepackage[margin=2.2cm]{geometry}
\usepackage{graphicx}
\usepackage{listings}
\usepackage{xcolor}
\usepackage{enumitem}
\usepackage{booktabs}
\usepackage{amsmath}
\usepackage[hidelinks]{hyperref}
\usepackage{fancyhdr}
\usepackage{float}

% ---------- 主题色（改成你喜欢的颜色） ----------
\definecolor{theme}{RGB}{40,80,140}

% ---------- 页眉页脚 ----------
\pagestyle{fancy}
\fancyhf{}
\fancyhead[L]{\small\color{theme}系统开发工具基础 · 实验报告}
\fancyhead[R]{\small 甘如嫣}
\fancyfoot[C]{\small\thepage}
\renewcommand{\headrule}{\color{theme}\rule{\headwidth}{0.4pt}}

% ---------- 代码块样式 ----------
\lstset{
  basicstyle=\ttfamily\small,
  breaklines=true,
  frame=single,
  backgroundcolor=\color{gray!8},
  keywordstyle=\color{blue!70!black},
  commentstyle=\color{green!45!black},
  showstringspaces=false,
  columns=fullflexible,
}

% ---------- 截图快捷命令 ----------
% \shot{图注}{图片相对路径}{宽度}
\newcommand{\shot}[3]{\begin{figure}[H]\centering\includegraphics[width=#3]{#2}\caption{#1}\end{figure}}

\begin{document}

% ============================================================
% 封面
% ============================================================
\begin{titlepage}
\centering
\vspace*{3cm}
{\Huge\bfseries\color{theme} 实验报告\par}
\vspace{0.8cm}
{\LARGE 系统开发工具基础（第三周）\par}
\vspace{2.5cm}
\begin{tabular}{rl}
\textbf{姓名} & 甘如嫣 \\
\textbf{学号} & 25090032014 \\
\textbf{日期} & \today \\
\end{tabular}
\vfill
\end{titlepage}

\tableofcontents
\newpage

% ============================================================
\section{GitHub 仓库与提交记录}

\subsection{仓库链接}
仓库地址：\url{https://github.com/ganruyan/system-basics-tools}

\subsection{提交记录}
本仓库按题、按实例增量提交（非单次全量提交）。本周（第三周）共提交 \textbf{18} 次 commit，
覆盖第 9～12 题、实例 1～11、解题感悟与实验报告，每次提交只针对一个逻辑单元，保留清晰的提交历史。

\shot{GitHub 提交记录（第三周增量提交）}{figs/commit_records.png}{0.92\textwidth}

% ============================================================
\clearpage
\section{课上给助教检查的内容与结果}

本节课布置并给助教检查的共 \textbf{4 道题目}（第 9～12 题），主题分别覆盖 Python 打包、
智能体编程、协作材料改写与 PyTorch 训练。下面每题先放题目原文，再放运行结果截图。

% ============ 题目 9 ============
\subsection{题目 9\quad 从源码构建并在干净环境安装 Wheel}

\textbf{主题：}Packaging and Shipping Code\quad\textbf{建议用时：}12—15 分钟

\textbf{题目原文：}在 q09 中按给定内容创建一个最小 Python 命令行包，并验证 wheel 安装。

\textbf{给定内容（\texttt{src/greetlab/cli.py}）：}
\begin{lstlisting}[language=Python]
import argparse

def main():
    p = argparse.ArgumentParser()
    p.add_argument("--name", required=True)
    a = p.parse_args()
    print(f"Hello, {a.name}!")
\end{lstlisting}

\textbf{给定内容（\texttt{pyproject.toml}）：}
\begin{lstlisting}
[build-system]
requires = ["setuptools>=68"]
build-backend = "setuptools.build_meta"
[project]
name = "greetlab-学号"
version = "0.1.0"
[project.scripts]
sdt-greet = "greetlab.cli:main"
\end{lstlisting}

\textbf{题目要求：}
\begin{itemize}
  \item 创建 \texttt{src/greetlab/\_\_init\_\_.py}、给定的 \texttt{cli.py} 和 \texttt{pyproject.toml}，并把「学号」替换为自己的学号。
  \item 在已安装 build 工具的课程环境中运行 \texttt{python -m build}，生成 wheel。
  \item 新建一个干净虚拟环境，只从生成的 wheel 安装。
  \item 切换到 q09 之外运行 \texttt{sdt-greet --name <学号>}，确认输出正确。
\end{itemize}

\textbf{运行结果：}
\shot{创建包文件：pyproject.toml 与 cli.py}{figs/q09_1.png}{0.45\textwidth}
\shot{greetlab 包目录结构}{figs/q09_2.png}{0.6\textwidth}
\shot{python -m build 成功生成 wheel}{figs/q09_3.png}{0.85\textwidth}
\shot{干净虚拟环境中 pip install 安装 wheel}{figs/q09_4.png}{0.6\textwidth}
\shot{运行 sdt-greet --name 25090032014，输出正确}{figs/q09_5.png}{0.85\textwidth}

% ============ 题目 10 ============
\clearpage
\subsection{题目 10\quad 让编程智能体进入可验证的修复循环}

\textbf{主题：}智能体编程\quad\textbf{建议用时：}12—15 分钟

\textbf{题目原文：}复用第 9 题的 greetlab 包，使用编程智能体完成一个小型测试驱动修复。

\textbf{题目要求：}
\begin{itemize}
  \item 复制 q09 为 q10，添加一个失败测试：当 name 只含空白字符时，main 应以 \texttt{SystemExit(2)} 结束。
  \item 先运行测试确认失败，再向编程智能体说明目标、约束和测试命令，要求它修改实现并运行测试。
  \item 人工检查 diff，撤销无关修改，并再次运行测试。
  \item 在 \texttt{ai\_log.md} 中用不超过 5 行记录核心提示、智能体改动和人工验证。
\end{itemize}

\textbf{运行结果：}
\shot{先运行测试，确认失败（FAILED）}{figs/q10_1.png}{0.85\textwidth}
\shot{向编程智能体说明目标与约束}{figs/q10_2.png}{0.5\textwidth}
\shot{人工检查 diff，确认仅改动校验逻辑}{figs/q10_3.png}{0.85\textwidth}
\shot{智能体修改后的 cli.py}{figs/q10_4.png}{0.85\textwidth}
\shot{ai\_log.md 记录核心提示、改动与验证}{figs/q10_5.png}{0.42\textwidth}
\shot{再次运行 pytest，1 passed}{figs/q10_6.png}{0.85\textwidth}

% ============ 题目 11 ============
\clearpage
\subsection{题目 11\quad 把「无法处理」的协作材料改成可执行信息}

\textbf{主题：}不止于代码\quad\textbf{建议用时：}12—15 分钟

\textbf{题目原文：}围绕「空姓名仍输出问候语」这一缺陷，改写下面三段质量较差的协作材料。

\textbf{给定内容：}
\begin{itemize}
  \item Issue：Windows 上运行不了，尽快修复。
  \item 提交信息：fix bug
  \item 评审意见：这里写得不好，重写。
  \item 已知事实：运行 \texttt{sdt-greet --name " "} 时仍输出 \texttt{Hello, !}，并以 0 退出。
\end{itemize}

\textbf{题目要求：}
\begin{itemize}
  \item 在 \texttt{communication.md} 中重写 Issue，包含环境、复现命令、期望结果和实际结果；未知信息写「待确认」。
  \item 重写提交信息：标题使用祈使语气，正文说明问题和解决方案。
  \item 重写评审意见：指出具体行为、风险和建议动作，并标注 Blocking、Suggestion 或 Nit。
  \item 全文控制在 400 字以内，保持专业、简洁。
\end{itemize}

\textbf{运行结果：}
\shot{communication.md 改写后的内容}{figs/q11_1.png}{0.85\textwidth}

% ============ 题目 12 ============
\clearpage
\subsection{题目 12\quad 修复一个可复现的线性回归训练循环}

\textbf{主题：}Python 与 PyTorch\quad\textbf{建议用时：}12—15 分钟

\textbf{题目原文：}将下方代码保存为 \texttt{q12/train.py}，补全 TODO 并完成一次 PyTorch CPU 训练。

\textbf{给定内容：}
\begin{lstlisting}[language=Python]
import torch
from torch import nn

torch.manual_seed(20260907)
x = torch.linspace(-1, 1, 101).reshape(-1, 1)
y = 3 * x - 1
model = nn.Linear(1, 1)
loss_fn = nn.MSELoss()
opt = torch.optim.SGD(model.parameters(), lr=0.1)

for _ in range(200):
    pred = model(x)
    loss = loss_fn(pred, y)
    # TODO：清空梯度、反向传播、更新参数

    # TODO：进入评估模式并在 no_grad 中打印最终损失、weight 和 bias
\end{lstlisting}

\textbf{题目要求：}
\begin{itemize}
  \item 补全训练循环，正确调用 \texttt{zero\_grad}、\texttt{backward} 和 \texttt{step}。
  \item 训练后调用 \texttt{model.eval()}，并在 \texttt{torch.no\_grad()} 中计算最终损失。
  \item 打印最终损失、weight 和 bias；最终损失应小于 0.001。
  \item 不得直接把 weight 和 bias 赋值为 3 和 $-1$。
\end{itemize}

\textbf{运行结果：}
\shot{train.py 补全后运行，输出 loss / weight / bias}{figs/q12_1.png}{0.8\textwidth}

% ============================================================
\clearpage
\section{课后练习（>10 个实例）}

本次课后练习共完成 \textbf{11 个实例}（实例 1～11），满足「每次课完成 >10 个实例」的要求。
下面每个实例先放题目原文，再放运行结果截图。

% ============ 实例 1 ============
\subsection{实例 1\quad printenv 对比虚拟环境前后差异}

\textbf{题目原文：}Save your environment with printenv to a file, create a venv, activate it, printenv to another file and diff before.txt after.txt. What changed in the environment? Why does the shell prefer the venv? (Hint: look at \texttt{\$PATH} before and after activation.) Run which deactivate and reason about what the deactivate bash function is doing.

\textbf{运行结果：}
\shot{printenv 导出环境变量，激活 venv 后 diff 前后差异}{figs/inst01.png}{0.55\textwidth}

% ============ 实例 2 ============
\clearpage
\subsection{实例 2\quad 创建 Python 包并生成 lockfile}

\textbf{题目原文：}Create a Python package with a pyproject.toml and install it in a virtual environment. Create a lockfile and inspect it.

\textbf{运行结果：}
\shot{创建 Python 包（pyproject.toml）}{figs/inst02_1.png}{0.85\textwidth}
\shot{在虚拟环境安装并生成 lockfile}{figs/inst02_2.png}{0.85\textwidth}

% ============ 实例 3 ============
\clearpage
\subsection{实例 3\quad 编写 Python 应用 Dockerfile}

\textbf{题目原文：}Write a Dockerfile for a simple Python application

\textbf{运行结果：}
\shot{编写的 Python 应用 Dockerfile}{figs/inst03.png}{0.85\textwidth}

% ============ 实例 4 ============
\clearpage
\subsection{实例 4\quad docker-compose 运行应用与 Redis 缓存}

\textbf{题目原文：}Then write a docker-compose.yml that runs your application alongside a Redis cache

\textbf{运行结果：}
\shot{docker-compose.yml（应用 + Redis 缓存）}{figs/inst04.png}{0.85\textwidth}

% ============ 实例 5 ============
\clearpage
\subsection{实例 5\quad 三种方式实现两数之和}

\textbf{题目原文：}把同一个简单编程任务（选取 LeetCode 两数之和）分别用手动编码、AI 自动补全、AI 聊天对话三种方式实现，对比三者体验差异；可选尝试智能体模式。

\textbf{运行结果：}
\shot{手动编码 / AI 补全 / AI 对话三种实现两数之和}{figs/inst05.png}{0.8\textwidth}

% ============ 实例 6 ============
\clearpage
\subsection{实例 6\quad 用 AI 生成命令处理 Markdown}

\textbf{题目原文：}利用 AI 完成 Markdown 无序列表提取任务。思考 AI 直接修改文件存在什么缺点。让 AI 输出 shell 命令（grep/sed），人工手动执行命令完成文本筛选，禁止 AI 直接修改原始 md 文件。

\textbf{运行结果：}
\shot{用 AI 生成 grep/sed 命令，人工执行处理 Markdown 无序列表}{figs/inst06.png}{0.85\textwidth}

% ============ 实例 7 ============
\clearpage
\subsection{实例 7\quad 氛围编程一个小应用}

\textbf{题目原文：}从零开始氛围编程一个小应用。不要手动编写一行代码。

\textbf{运行结果：}
\shot{氛围编程生成的小应用运行结果}{figs/inst07.png}{0.5\textwidth}

% ============ 实例 8 ============
\clearpage
\subsection{实例 8\quad 原地修改与输出新文件对比}

\textbf{题目原文：}对比两种文本处理思路：直接原地修改源文件 和 先输出结果到新文件、不改动原始文件，分别简述优缺点，并说明日常处理重要文本时优先选择哪种方案。

\textbf{运行结果：}
\shot{原地修改与输出新文件两种思路对比}{figs/inst08.png}{0.85\textwidth}

% ============ 实例 9 ============
\clearpage
\subsection{实例 9\quad YOLO 模式风险与隔离方案}

\textbf{题目原文：}调研 AI 智能体 YOLO 自主模式的含义，分析该模式存在的安全风险，写出一份基于 Ubuntu 虚拟机的隔离运行方案，不需要真实搭建部署。

\textbf{运行结果：}
\shot{YOLO 自主模式含义、风险与隔离运行方案}{figs/inst09.png}{0.7\textwidth}

% ============ 实例 10 ============
\clearpage
\subsection{实例 10\quad 查找 .py 文件并遍历脚本}

\textbf{题目原文：}在家目录下，编写一条命令完成：查找当前目录下所有 \texttt{.py} 后缀文件，打印文件名称、文件权限；随后写一段简单 bash 脚本，遍历目录输出所有 Python 文件。

\textbf{运行结果：}
\shot{查找 .py 文件（名称+权限）与遍历脚本输出}{figs/inst10.png}{0.85\textwidth}

% ============ 实例 11 ============
\clearpage
\subsection{实例 11\quad Markdown 无序列表文本流水线}

\textbf{题目原文：}对 Markdown 文件进行文本流水线处理：提取无序列表、去除空白行、把开头 \texttt{- } 统一替换成 \texttt{* }，最终保存到新文件，原始文件保持不变。

\textbf{运行结果：}
\shot{Markdown 无序列表文本流水线处理结果}{figs/inst11.png}{0.85\textwidth}

% ============================================================
\clearpage
\section{解题感悟}

这次完成实验题目，接触了 Docker 配置、Python 代码和 Linux 命令。一开始我只是照着步骤敲命令，很多地方没有深入思考。最开始只用简单 grep 提取文本，做完觉得功能很单薄，后来尝试管道组合多条命令，才感受到 Shell 处理文本的方便。

手动写代码和借助工具生成代码对比下来，我发现工具虽然省事，但生成的内容需要自己检查，不能直接照搬。写配置文件的时候也注意到缩进、格式这类小细节很容易出错。实验过程里印象很深的一点，就是尽量不去改动原始文件，避免误操作弄丢数据。

整体做完收获挺多，不少命令和脚本还不太熟练，之后还要多动手练习，熟悉如何组合工具解决问题。

\end{document}
